% 中文摘要
\thispagestyle{empty}
\begin{center}
  {\CJKfamily{heiti}\zihao{-2}\bfseries 摘\quad 要}
\end{center}

\vspace{1em}

% 摘要正文（v3：贡献定位 + 真实数字）
随着大语言模型与工具调用技术的发展，AI Agent 已从单轮问答工具演进为能够理解目标、分解任务、调用外部工具并持续反馈的自治执行单元。多 Agent 协作时，系统需同时处理任务编排、跨进程通讯、状态一致性、执行隔离、实时观测与故障恢复等工程问题。现有多智能体框架以 Python 平台与角色对话流程为主，在本地化、高性能、可审计的分布式通讯底座方面仍有较大空间；尤其是 Agent2Agent (A2A) 协议虽已有官方 Rust SDK（a2aproject 维护的 a2a-rs）与若干第三方 crate，但缺少与本地优先平台的事件总线、编排层与人工介入语义深度集成的实现。

针对上述问题，本文设计并实现一套名为 Fuxi 的 Rust 平台。系统采用「玄女—门客」分层协作模式将面向用户的顶层 Agent 与执行任务的工作 Agent 解耦；自主从零实现一份面向 Fuxi 场景的 A2A 风格核心 RPC 闭环（覆盖 agent discovery、task send、task stream、task query 与 task cancel 五条主路径，沿用早期 A2A JSON-RPC binding），并把 A2A 规范已有的 input-required 中间状态映射到 Fuxi 内部 \texttt{Task::PendingApproval} 与 \texttt{ShelfStatus::AwaitingInput}，使人工介入成为平台级可观测事件；以 Tokio broadcast 与 SQLite WAL 构建实时事件总线，使实时推送与历史回放在同一抽象上共存；以 Firehose、WebSocket、SSE 与 IM API 提供多端观测；以 Git worktree 与三层 sandbox 保证任务执行隔离。论文进一步给出端到端延迟分解、广播事件延迟下界与 Worker 选择函数的形式化定义，使系统性能具备可量化、可比较的分析基础。

实验环节在 Apple Silicon 平台上对吞吐量、延迟、参数敏感性与事件总线压力进行了多维测量。结果表明，在剥离 LLM 推理后的合成任务基准下（每任务以 10 ms sleep 模拟），Fuxi 在 8 worker 配置下达到 665.00 tasks/s 的调度路径吞吐，扩展至 16 worker 时增长至 1288.24 tasks/s（1$\rightarrow$16 worker 区间扩展效率介于 80.5\%--83.1\%，dispatcher 在该区间未饱和）；任务派发链路 p50 延迟为 32.93 ms（含 HTTP / SQLite / 反序列化等组件），进程内事件总线广播 p99 延迟为 0.26 ms；事件总线在 64 subscriber 并发订阅、10\,000 events/s 持续负载下保持零丢帧。需要说明的是，上述合成基准衡量的是编排与通讯层在剥离 LLM 推理后的纯通讯吞吐上限，真实 LLM Agent 任务的端到端吞吐受单次推理（典型 1--30 s）主导。这些数据从侧面验证了「事件驱动与追加式日志结合的通讯架构能够在本地优先的 AI Agent 协作场景中同时兼顾吞吐、实时性与可追溯性」的设计假设，为后续多 Agent 系统的工程化提供了一套可借鉴的基础设施。

\vspace{1.5em}

\noindent {\CJKfamily{heiti}\bfseries 关键词}：AI Agent；分布式通讯；多智能体系统；事件总线；A2A 协议
